\documentclass[12pt]{article}
\usepackage[utf8]{inputenc}
\usepackage[T1]{fontenc}
\usepackage[margin=1in]{geometry}
\usepackage{setspace}
\usepackage{xcolor}
\usepackage[colorlinks=true, urlcolor=blue!60!black, linkcolor=blue!60!black]{hyperref}
\setlength{\parindent}{2em}
\setlength{\parskip}{1.25em}
\renewcommand{\baselinestretch}{1.0}
\begin{document}
\begin{center}
{\LARGE \textit{Duncan v. Louisiana}, 391 U.S. 145 (1968)\par}
\vspace{0.8em}
{\large Case Brief\par}
\end{center}
\vspace{1em}

\section{Memory Jogger}

In \href{https://www.courtlistener.com/opinion/107685/duncan-v-louisiana/}{\textit{Duncan v. Louisiana}}, the Supreme Court incorporated the Sixth Amendment right to a criminal jury trial against the states through the Fourteenth Amendment's Due Process Clause, holding that any state prosecution for a serious offense---including a misdemeanor punishable by up to two years' imprisonment---must afford the accused a trial by jury.

\section{Detailed Case Facts}

On October 18, 1966, Gary Duncan, a nineteen-year-old Black man, was driving along Highway 23 in Plaquemines Parish, Louisiana, when he saw two younger Black cousins standing by the side of the road in an encounter with four white boys (\href{https://supreme.justia.com/cases/federal/us/391/145/}{391 U.S. at 147}). Knowing that his cousins had recently transferred to a formerly all-white public high school and had reported racial incidents there, Duncan stopped his car, got out, and urged his cousins to break off the encounter and get into his vehicle. Witnesses agreed that Duncan spoke with the white boys and was about to enter his car to drive away, but the testimony sharply diverged on what happened next: the white boys and a white onlooker testified that Duncan slapped one of the white boys, Herman Landry, on the elbow, whereas Duncan and his cousins testified that Duncan had merely touched Landry's elbow without slapping him (\href{https://caselaw.findlaw.com/court/us-supreme-court/391/145.html}{391 U.S. at 147}).

Duncan was arrested and charged with simple battery under Louisiana Revised Statutes \S~14:35 (1950), which defined simple battery as ``a battery, without the consent of the victim, committed without a dangerous weapon,'' and classified it as a misdemeanor punishable by a fine of not more than \$300, imprisonment for not more than two years, or both (\href{https://supreme.justia.com/cases/federal/us/391/145/}{391 U.S. at 146 \& n.1}). Under Article VII, \S~41 of the Louisiana Constitution, however, jury trials were granted only in cases in which capital punishment or imprisonment at hard labor could be imposed; all offenses not punishable at hard labor were tried to a judge alone. Because simple battery did not carry hard labor, the trial judge denied Duncan's timely request for a jury trial. Following a bench trial, the judge found Duncan guilty beyond a reasonable doubt and sentenced him to 60 days in the parish prison and a \$150 fine (\href{https://caselaw.findlaw.com/court/us-supreme-court/391/145.html}{391 U.S. at 146}).

The prosecution unfolded against the volatile backdrop of federally mandated school desegregation in Plaquemines Parish, a political stronghold dominated by staunch segregationist and former District Attorney Leander Perez, whose son Leander H. Perez, Jr. appeared on Louisiana's Supreme Court brief (\href{https://law.stanford.edu/2022/10/12/the-civil-rights-significance-of-duncan-v-louisiana-and-the-right-to-a-jury-trial/}{Stanford Center for Racial Justice}; \href{https://caselaw.findlaw.com/court/us-supreme-court/391/145.html}{FindLaw}). Civil rights lawyers Richard B. Sobol, Alvin J. Bronstein, and Anthony G. Amsterdam represented Duncan, viewing Louisiana's denial of a jury trial in misdemeanor prosecutions as a tool used to harass Black citizens and chill participation in school integration. Constitutionally, the case squarely tested whether the Sixth Amendment guarantee---``In all criminal prosecutions, the accused shall enjoy the right to a speedy and public trial, by an impartial jury''---applies to the states via the Due Process Clause of the Fourteenth Amendment.

\section{Procedural History}

\begin{itemize}
  \item \textbf{Trial Court:} In the Twenty-fifth Judicial District Court of Louisiana (Plaquemines Parish), Gary Duncan demanded a jury trial on the charge of simple battery under La. Rev. Stat. \S~14:35. Relying on La. Const. Art. VII, \S~41, the trial judge denied the request, tried Duncan without a jury, convicted him, and sentenced him to 60 days in the parish prison and a \$150 fine (\href{https://caselaw.findlaw.com/court/us-supreme-court/391/145.html}{391 U.S. at 146}).
  \item \textbf{Louisiana Supreme Court:} Duncan sought review in the Supreme Court of Louisiana, arguing that denying him a jury trial violated the Sixth and Fourteenth Amendments to the U.S. Constitution. Finding ``[n]o error of law in the ruling complained of,'' the Louisiana Supreme Court denied a writ of certiorari, 250 La. 253, 195 So. 2d 142 (1967) (\href{https://supreme.justia.com/cases/federal/us/391/145/}{391 U.S. at 146 \& n.2}).
  \item \textbf{U.S. Supreme Court:} Duncan appealed as of right pursuant to 28 U.S.C. \S~1257(2). The Supreme Court noted probable jurisdiction, 389 U.S. 809 (1967), and set the case for oral argument alongside \href{https://supreme.justia.com/cases/federal/us/391/194/}{\textit{Bloom v. Illinois}}, 391 U.S. 194 (1968), on January 17, 1968. On May 20, 1968, the Court reversed the judgment below and remanded for further proceedings (\href{https://www.courtlistener.com/opinion/107685/duncan-v-louisiana/}{391 U.S. at 162}).
\end{itemize}

\section{Judicial Votes}

The Supreme Court decided \textit{Duncan v. Louisiana} by a \textbf{7--2} vote:

\begin{itemize}
  \item \textbf{Majority (7):} Justice Byron R. White delivered the opinion of the Court, joined by Chief Justice Earl Warren and Justices Hugo L. Black, William O. Douglas, William J. Brennan, Jr., Abe Fortas, and Thurgood Marshall (\href{https://supreme.justia.com/cases/federal/us/391/145/}{391 U.S. at 146--162}).
  \item \textbf{Concurrence (Black):} Justice Hugo L. Black filed a concurring opinion, joined by Justice William O. Douglas (\href{https://supreme.justia.com/cases/federal/us/391/145/}{391 U.S. at 162--171}).
  \item \textbf{Concurrence (Fortas):} Justice Abe Fortas filed a separate concurring opinion applicable to both \textit{Duncan} and \textit{Bloom v. Illinois} (\href{https://supreme.justia.com/cases/federal/us/391/194/}{391 U.S. at 211--215}).
  \item \textbf{Dissent (2):} Justice John Marshall Harlan II filed a dissenting opinion, joined by Justice Potter Stewart (\href{https://supreme.justia.com/cases/federal/us/391/145/}{391 U.S. at 171--193}).
\end{itemize}

\section{Holding}

\textbf{Judgment:} Reversed and remanded. Louisiana violated the Fourteenth Amendment by denying Gary Duncan a jury trial on a charge of simple battery punishable by up to two years' imprisonment.

\textbf{Rule Announced:} Because trial by jury in criminal cases is fundamental to the American scheme of justice, the Fourteenth Amendment's Due Process Clause guarantees a right to a jury trial in all state criminal prosecutions that, if tried in a federal court, would fall within the Sixth Amendment's guarantee (\href{https://supreme.justia.com/cases/federal/us/391/145/}{391 U.S. at 149}). While ``petty offenses'' remain exempt from the jury-trial requirement, the maximum penalty authorized by the legislature---not merely the sentence actually imposed---is of major relevance in determining whether a crime is serious, and any crime carrying a maximum authorized penalty of two years in prison is a serious crime entitling the defendant to a jury trial (\href{https://supreme.justia.com/cases/federal/us/391/145/}{391 U.S. at 159--162}).

\textbf{Holding vs. Dicta:} The Court's binding holdings are (1) that the Sixth Amendment right to a criminal jury trial is incorporated against the states through the Fourteenth Amendment, and (2) that an offense carrying a statutory maximum of two years' imprisonment is a serious crime requiring a jury trial regardless of the lighter 60-day sentence actually imposed. By contrast, the Court expressly declined to fix the exact boundary between petty and serious offenses, leaving that line for future cases (later drawn at six months' authorized imprisonment in \textit{Baldwin v. New York}, 399 U.S. 66 (1970)). Likewise, Footnote 30's remarks that the Court's prior federal Sixth Amendment decisions on twelve-person juries and unanimous verdicts were ``always subject to reconsideration'' and unlikely to require widespread state changes (\href{https://supreme.justia.com/cases/federal/us/391/145/}{391 U.S. at 158--159 n.30}) were dicta that foreshadowed decades of litigation over jury size and unanimity.

\section{Analysis of Opinions}

\subsection{Majority Opinion (Justice White)}

Justice White reframed the Fourteenth Amendment selective-incorporation inquiry in a pivotal methodological shift. Rather than asking, as \textit{Palko v. Connecticut}, 302 U.S. 319 (1937), had, whether a hypothetical ``fair and enlightened system of justice'' could exist anywhere in the world without a jury, the Court asked whether a procedural safeguard is ``fundamental to the American scheme of justice''---that is, ``necessary to an Anglo-American regime of ordered liberty'' (\href{https://supreme.justia.com/cases/federal/us/391/145/}{391 U.S. at 149--150 \& n.14}). While a fair nonjury system might be imagined in the abstract, every American state operated within a common-law framework built around the criminal jury.

To establish the jury's fundamental status, White surveyed its centuries-long pedigree from Magna Carta and the English Bill of Rights of 1689, through Blackstone's praise of the jury as a ``strong and two-fold barrier'' against Crown prerogative, to the Stamp Act Congress (1765), the First Continental Congress (1774), the Declaration of Independence, Article III, \S~2, the Sixth Amendment, and the constitutions of all fifty states (\href{https://supreme.justia.com/cases/federal/us/391/145/}{391 U.S. at 151--154}). Functionally, the criminal jury protects the accused against ``oppression by the Government''---specifically, ``against the corrupt or overzealous prosecutor and against the compliant, biased, or eccentric judge''---by insisting on community participation before the state deprives a citizen of liberty (\href{https://supreme.justia.com/cases/federal/us/391/145/}{391 U.S. at 155--156}). White rejected contrary language in \textit{Maxwell v. Dow}, 176 U.S. 581 (1900), \textit{Snyder v. Massachusetts}, 291 U.S. 97 (1934), and \textit{Palko} as nonbinding dicta or discredited doctrine, and cited Harry Kalven, Jr. and Hans Zeisel's empirical study, \textit{The American Jury} (1966), to refute claims that lay jurors are incompetent factfinders (\href{https://supreme.justia.com/cases/federal/us/391/145/}{391 U.S. at 154--157}).

Turning to Louisiana's fallback argument that Duncan's 60-day sentence made his offense ``petty,'' White acknowledged the historical petty-offense exception but held that the statutory maximum penalty---here, two years---serves as the primary objective gauge of society's judgment of an offense's gravity, distinguishing \textit{Cheff v. Schnackenberg}, 384 U.S. 373 (1966), as a criminal contempt case lacking a statutory maximum (\href{https://supreme.justia.com/cases/federal/us/391/145/}{391 U.S. at 159--162 \& n.35}). Noting that the federal system drew the petty-offense line at six months and that 49 of the 50 states limited nonjury trials to crimes punishable by no more than one year, White held that a crime punishable by two years' imprisonment is plainly serious (\href{https://supreme.justia.com/cases/federal/us/391/145/}{391 U.S. at 161--162}).

\subsection{Concurring Opinion (Justice Black, joined by Justice Douglas)}

Justice Black joined the Court's opinion while reaffirming his celebrated ``total incorporation'' thesis from his dissent in \textit{Adamson v. California}, 332 U.S. 46 (1947): that Section 1 of the Fourteenth Amendment---especially its Privileges or Immunities Clause---was historically intended by its congressional sponsors, Representative John Bingham and Senator Jacob Howard, to apply the entire first eight amendments directly against the states (\href{https://supreme.justia.com/cases/federal/us/391/145/}{391 U.S. at 162--167}). Black forcefully rebutted Professor Charles Fairman's 1949 \textit{Stanford Law Review} critique, arguing that Fairman relied improperly on silence in state ratification debates rather than the explicit floor statements of the Amendment's sponsors (\href{https://supreme.justia.com/cases/federal/us/391/145/}{391 U.S. at 165--167}).

Black's deepest objection to Justice Harlan's ``fundamental fairness'' approach was anti-discretionary: treating due process as an ``evolving concept'' untethered to the Bill of Rights' text gave judges ``unconfined power'' to strike down or uphold state laws according to personal ethics (\href{https://supreme.justia.com/cases/federal/us/391/145/}{391 U.S. at 168--170}). Selective incorporation, though less historically pure than total incorporation, was vastly preferable because it anchored judicial review to specific constitutional guarantees and had already extended nearly all of the Bill of Rights to the states (\href{https://supreme.justia.com/cases/federal/us/391/145/}{391 U.S. at 171}).

\subsection{Concurring Opinion (Justice Fortas)}

Writing in an opinion covering both \textit{Duncan} and \textit{Bloom v. Illinois}, Justice Fortas joined the majority's holding that due process requires a jury trial in serious state prosecutions, agreeing that the Court's understanding of ``due process of law'' deepens over time alongside evolving conceptions of human dignity (\href{https://supreme.justia.com/cases/federal/us/391/194/}{391 U.S. at 211--212}). Crucially, however, Fortas rejected ``jot-for-jot'' incorporation---the assumption that ``the tail must go with the hide''---arguing that incorporating the core right to a jury trial should not automatically saddle the states with every ancillary federal rule, such as twelve-person juries or unanimous verdicts (\href{https://supreme.justia.com/cases/federal/us/391/194/}{391 U.S. at 213}). Emphasizing that ``[o]ur Constitution sets up a federal union, not a monolith,'' Fortas urged the Court to welcome state procedural variations that preserve basic fairness without impairing the core function of trial by jury (\href{https://supreme.justia.com/cases/federal/us/391/194/}{391 U.S. at 214--215}).

\subsection{Dissenting Opinion (Justice Harlan, joined by Justice Stewart)}

Justice Harlan mounted a sweeping structural and historical attack on selective incorporation, calling the Court's approach ``an uneasy and illogical compromise'' that enjoyed neither the historical consistency of Black's total incorporation nor the principled discrimination of traditional due process analysis (\href{https://supreme.justia.com/cases/federal/us/391/145/}{391 U.S. at 172, 176, 180--181}). Endorsing Professor Fairman's historical scholarship, Harlan contended that the Fourteenth Amendment's framers never intended to incorporate the Bill of Rights ``jot-for-jot and case-for-case,'' but instead wrote broad terms requiring a ``gradual process of judicial inclusion and exclusion'' focused on whether a state procedure was ``fundamentally fair'' (\href{https://supreme.justia.com/cases/federal/us/391/145/}{391 U.S. at 174--181}).

On the merits, Harlan argued that bench trials are not inherently unfair---as proven by the prevalence of jury waivers in the United States and summary trials in England---and that the Court itself conceded Duncan's bench trial was not shown to be unfair (\href{https://supreme.justia.com/cases/federal/us/391/145/}{391 U.S. at 187--190}). He warned that forcing federal Sixth Amendment rules onto all fifty states would either put states in a ``constitutional straitjacket'' or tempt the Court to ``water[] down'' federal Bill of Rights protections in the pursuit of uniformity (\href{https://supreme.justia.com/cases/federal/us/391/145/}{391 U.S. at 175--176, 182 n.21}). Invoking Justice Brandeis's famous ``laboratory'' metaphor from \textit{New State Ice Co. v. Liebmann}, 285 U.S. 262 (1932), Harlan concluded that states should remain free to calibrate where the line between jury and nonjury trials belongs (\href{https://supreme.justia.com/cases/federal/us/391/145/}{391 U.S. at 192--193}).

\section{Examples: Future Applications}

\begin{enumerate}
  \item \textbf{Same-Side (Jury Required) --- State Misdemeanor Carrying One Year in Jail:} A state statute punishes first-offense shoplifting of goods valued over \$500 by up to twelve months in county jail, but denies a jury trial in municipal court. Under \textit{Duncan}, the authorized maximum penalty---not the prosecutor's promise to seek only a 30-day sentence---governs whether the offense is serious; because twelve months exceeds the traditional six-month petty-offense benchmark cited by the Court, the defendant is constitutionally entitled to a jury trial.
  \item \textbf{Same-Side (Jury Required) --- Serious Criminal Contempt Without a Statutory Cap:} A state trial judge holds a defense lawyer in criminal contempt for violating a gag order and, under a state rule setting no statutory maximum penalty, summarily sentences the lawyer to eighteen months in prison without a jury. Applying \textit{Duncan} and its companion case \href{https://supreme.justia.com/cases/federal/us/391/194/}{\textit{Bloom v. Illinois}}, when the legislature has not fixed a maximum sentence, the penalty actually imposed serves as the best measure of seriousness; an eighteen-month sentence makes the contempt a serious crime requiring a jury trial.
  \item \textbf{Opposite-Side (No Jury Required) --- Petty Traffic Violation:} A motorist is charged with reckless driving under a state statute carrying a maximum penalty of 90 days in jail and a \$250 fine, and the municipal court denies a jury trial. Under \textit{Duncan}'s explicit preservation of the petty-offense exception (\href{https://supreme.justia.com/cases/federal/us/391/145/}{391 U.S. at 159--161}), an offense carrying no more than six months' authorized imprisonment and modest collateral penalties falls within the petty-offense category and may be tried to a judge alone.
  \item \textbf{Opposite-Side (No Jury Required) --- Civil Regulatory Penalty:} A state environmental agency initiates a civil enforcement action seeking \$50,000 in statutory civil forfeitures against a factory owner for unlawful wastewater discharge, with no possibility of incarceration. Because \textit{Duncan} incorporated the Sixth Amendment right to a jury trial ``in criminal cases'' (\href{https://supreme.justia.com/cases/federal/us/391/145/}{391 U.S. at 149}), and the Seventh Amendment civil jury right remains unincorporated against the states, \textit{Duncan} does not mandate a jury trial in noncriminal regulatory proceedings.
  \item \textbf{Fence-Sitter --- Stacked Petty Charges Carrying Aggregate Exposure Over Six Months:} A state prosecutor charges a protester in a single bench trial with four separate counts of disorderly conduct, each carrying a maximum penalty of 90 days in jail (an aggregate exposure of 360 days). On one hand, \textit{Duncan}'s focus on the defendant's total exposure to serious deprivation of liberty suggests a jury should be required when consecutive sentences could reach a year; on the other hand, \textit{Duncan} treats the legislature's penalty for ``a particular crime'' as the objective gauge of seriousness (\href{https://supreme.justia.com/cases/federal/us/391/145/}{391 U.S. at 159--160}), a tension the Supreme Court later resolved against aggregation in \textit{Lewis v. United States}, 518 U.S. 322 (1996).
\end{enumerate}

\section{Critique}

\textit{Duncan v. Louisiana} stands as a watershed of the Warren Court's criminal procedure revolution, yet it has drawn sustained scrutiny across doctrinal and ideological lines. Logically, Justice Harlan exposed a genuine tension at the heart of selective incorporation: by rejecting both total incorporation and pure case-by-case fairness, the majority offered no airtight textual explanation for why the Sixth Amendment jury right is ``in'' while the Fifth Amendment grand jury indictment right and Seventh Amendment civil jury right remain ``out.'' Moreover, Harlan's warning about ``watering down'' Bill of Rights protections proved prescient in the short term: in \textit{Williams v. Florida}, 399 U.S. 78 (1970), and \textit{Apodaca v. Oregon}, 406 U.S. 404 (1972), the Court diluted traditional common-law features of the jury (twelve jurors and unanimity) to accommodate state practices---a compromise the Court ultimately had to repudiate as to unanimity five decades later in \href{https://www.supremecourt.gov/opinions/19pdf/18-5924_n6io.pdf}{\textit{Ramos v. Louisiana}}, 590 U.S. 83 (2020).

From a \textbf{progressive and civil rights perspective}, \textit{Duncan} is celebrated for curbing arbitrary local justice in the Jim Crow South, yet scholars and Gary Duncan himself have noted its stark practical limits. Incorporating the jury right did not by itself guarantee an impartial or racially representative jury: Louisiana did not require juries to reflect a fair cross-section of the community until \textit{Taylor v. Louisiana}, 419 U.S. 522 (1975), and continued to permit nonunanimous convictions---a practice rooted in white supremacy---until \textit{Ramos} in 2020 (\href{https://law.stanford.edu/2022/10/12/the-civil-rights-significance-of-duncan-v-louisiana-and-the-right-to-a-jury-trial/}{Stanford Center for Racial Justice}). Indeed, Duncan later observed that had his case not been dismissed on remand after federal courts enjoined his bad-faith retrial in \textit{Duncan v. Perez}, 321 F. Supp. 181 (E.D. La. 1970) (\href{https://digitalcommons.mainelaw.maine.edu/cgi/viewcontent.cgi?article=1572&context=mlr}{Maine Law Review}), an all-white Plaquemines Parish jury would almost certainly have convicted him. Critics also note that \textit{Duncan}'s promise has been hollowed out by plea bargaining and by \textit{McKeiver v. Pennsylvania}, 403 U.S. 528 (1971), which withheld the jury right from juvenile delinquency proceedings.

From an \textbf{originalist and textualist perspective}, \textit{Duncan} elicits a fascinating split. Many modern originalists agree with Justice Black's ultimate conclusion that the Fourteenth Amendment incorporates fundamental Bill of Rights protections---including criminal jury trials---against the states, but argue that the majority chose the wrong textual vehicle: the Privileges or Immunities Clause, not the Due Process Clause, was the historically intended conduit for incorporation, as Justice Thomas later argued in \textit{McDonald v. City of Chicago}, 561 U.S. 742 (2010), and \textit{Ramos}. At the same time, originalists have criticized \textit{Duncan}'s post-1968 functionalist progeny (\textit{Williams}) for abandoning the Founding-era twelve-person jury requirement, arguing that an originalist reading of the Sixth Amendment actually protects criminal defendants more robustly than functional balancing (\href{https://scholarlycommons.law.northwestern.edu/cgi/viewcontent.cgi?article=7801&context=jclc}{Northwestern Journal of Criminal Law \& Criminology}). Conversely, federalism-minded conservatives echo Harlan's concern that selective incorporation via substantive due process stretches the Fourteenth Amendment's text beyond its 1868 meaning and displaces state experimentation in criminal administration.

\section{Key Quotations}

\begin{quote}
``Because we believe that trial by jury in criminal cases is fundamental to the American scheme of justice, we hold that the Fourteenth Amendment guarantees a right of jury trial in all criminal cases which --- were they to be tried in a federal court --- would come within the Sixth Amendment's guarantee.''
--- Justice White, Opinion of the Court (\href{https://supreme.justia.com/cases/federal/us/391/145/}{391 U.S. at 149}).
\end{quote}

\begin{quote}
``The question thus is whether given this kind of system a particular procedure is fundamental --- whether, that is, a procedure is necessary to an Anglo-American regime of ordered liberty.''
--- Justice White, Opinion of the Court (\href{https://supreme.justia.com/cases/federal/us/391/145/}{391 U.S. at 149--150 n.14}).
\end{quote}

\begin{quote}
``Providing an accused with the right to be tried by a jury of his peers gave him an inestimable safeguard against the corrupt or overzealous prosecutor and against the compliant, biased, or eccentric judge. If the defendant preferred the common-sense judgment of a jury to the more tutored but perhaps less sympathetic reaction of the single judge, he was to have it.''
--- Justice White, Opinion of the Court (\href{https://caselaw.findlaw.com/court/us-supreme-court/391/145.html}{391 U.S. at 156}).
\end{quote}

\begin{quote}
``What more precious `privilege' of American citizenship could there be than that privilege to claim the protections of our great Bill of Rights?''
--- Justice Black, Concurring (\href{https://supreme.justia.com/cases/federal/us/391/145/}{391 U.S. at 166}).
\end{quote}

\begin{quote}
``The Due Process Clause commands us to apply its great standard to state court proceedings to assure basic fairness. It does not command us rigidly and arbitrarily to impose the exact pattern of federal proceedings upon the 50 States. ... Our Constitution sets up a federal union, not a monolith.''
--- Justice Fortas, Concurring (\href{https://supreme.justia.com/cases/federal/us/391/194/}{391 U.S. at 213--214}).
\end{quote}

\begin{quote}
``In short, neither history nor sense supports using the Fourteenth Amendment to put the States in a constitutional straitjacket with respect to their own development in the administration of criminal or civil law.''
--- Justice Harlan, Dissenting (\href{https://supreme.justia.com/cases/federal/us/391/145/}{391 U.S. at 175--176}).
\end{quote}

\end{document}
